\section{Cancellation for an invertible differential}\label{sec:invertible}

We prove Theorem~\ref{thm:introcancellation}. The surviving cohomology
is $W$ in degree zero. Put $a=\arg z_1\in(0,\pi)$ and, in moving
orthonormal frames, define
\begin{equation}\label{eq:invertiblevariables}
 X_j=A_jD^{-1},\qquad H_X=\sum_jX_j^\dagger X_j,\qquad
 s=\tr H_X,\qquad \lambda=\lambda_{\min}(D^\dagger D).
\end{equation}

\begin{lemma}\label{lem:Xcontract}
Along the phase flow,
\begin{equation}\label{eq:Xcontract}
 \dot X_j=\tfrac12(X_j\Phi_U-\Phi_WX_j),\qquad
 s(t)\leq s(0)e^{-at}.
\end{equation}
\end{lemma}

\begin{proof}
Differentiate $A_jD^{-1}$ using \eqref{eq:mapflow}; the source phase
cancels. Since $A_j=X_jD$,
\begin{equation}\label{eq:targetX}
 \Omega_W=z_1\id-\sum_jX_j(\id-e^{-L})X_j^\dagger .
\end{equation}
The subtracted matrix is positive semidefinite. This target block
is normal, with scalar strictly positive imaginary part, and
$\arg(z_1-r)\geq\arg z_1$ for $r\geq0$. Thus $\Phi_W\geq a\id$,
whereas $\Phi_U<0$ on the degree-one branch. Consequently
\[
 \dot s=\sum_j\tr(X_j^\dagger X_j\Phi_U)
          -\tr\left(\sum_jX_jX_j^\dagger\Phi_W\right)
 \leq-as.
\]
\end{proof}

\begin{lemma}\label{lem:heatbounds}
Let $\kappa=\pi/(2e)$. For an invertible differential of rank $n$,
the coefficient terms in \eqref{eq:coeff} satisfy
\begin{equation}\label{eq:heatbounds}
 \norm{K_B(C)}_1\leq n\kappa s,\qquad
 \norm{J_L(E)}_1\leq ns,\qquad
 \norm{\Omega_W-z_1\id}_1\leq s.
\end{equation}
The constants do not depend on any singular-value ratio of $D$.
\end{lemma}

\begin{proof}
Write $D=U_DB^{1/2}$ and $S=U_D^\dagger H_XU_D$. Then
$C=B^{1/2}SB^{1/2}$, $S\geq0$ and $\tr S=s$.
For $u>0$,
$\norm{B^{1/2}e^{-uB}}\leq(2eu)^{-1/2}$.
Integration in \eqref{eq:K}, using
$\int_0^1[s(1-s)]^{-1/2}ds=\pi$, gives
$\norm{K_B(C)}\leq\kappa s$, hence the first trace-norm bound.

On $U$, $E=LH_XL$. In an eigenbasis of $L$ with eigenvalues $l_i$,
\[
 (J_L(E))_{ij}=k_{ij}(H_X)_{ij},\qquad
 k_{ij}=l_il_j\int_{\substack{x,z\geq0\\x+z\leq1}}
                  e^{-xl_i-zl_j}\,dx\,dz .
\]
Enlarging the integration region to the nonnegative quadrant gives
$0\leq k_{ij}\leq1$. Positivity of $H_X$ implies
$|(H_X)_{ij}|\leq\sqrt{(H_X)_{ii}(H_X)_{jj}}$. Therefore
\[
 |\inner{v}{J_L(E)v}|
 \leq\left(\sum_i|v_i|\sqrt{(H_X)_{ii}}\right)^2
 \leq s\norm v^2.
\]
Since $J_L(E)$ is Hermitian, this proves the second estimate.
Finally the positive semidefinite matrix subtracted in
\eqref{eq:targetX} has trace at most $s$.
\end{proof}

\begin{lemma}\label{lem:growinggap}
There is a constant $C$ depending on the initial data and charges such
that $\lambda(t)\geq at-C$ for every $t\geq0$.
\end{lemma}

\begin{proof}
The least eigenvalue $\lambda$ is locally Lipschitz. At almost every
time, choose a unit minimum eigenvector $v$ realizing its derivative,
and set $u=Dv/\sqrt\lambda$. Then $Lu=\lambda u$ and
\begin{equation}\label{eq:lambdaderivative}
 \dot\lambda
 =\lambda\bigl(\inner{v}{\Phi_Vv}-\inner{u}{\Phi_Uu}\bigr)>0.
\end{equation}
The same lower estimates hold on every vector in a repeated minimum
eigenspace, and hence for the corresponding one-sided eigenvalue
derivatives. In particular $\lambda\geq\lambda_0:=\lambda(0)>0$.

Apply \eqref{eq:weightedscalar} to
$\Omega_V=z_0e^{-B}+K_B(C)$ and
$-\Omega_U=z_0e^{-L}-J_L(E)$.
On $v$ and $u$ the reference heat eigenvalue is $e^{-\lambda}$.
The reference phases differ by $\pi$ after the degree-one branch
shift. Lemmas~\ref{lem:Xcontract} and \ref{lem:heatbounds} yield
\begin{equation}\label{eq:gapineq}
 \dot\lambda\geq\lambda
   \left[\pi-C_0e^{(\lambda-at)/2}\right],\qquad
 C_0=\pi\sqrt{\frac{ns(0)}{|z_0|}}(\sqrt\kappa+1).
\end{equation}
Let $L_0=\max\{0,2\log(2C_0/\pi)\}$, taking $L_0=0$ when $C_0=0$.
Whenever $\lambda\leq at-L_0$, the bracket in \eqref{eq:gapineq}
is at least $\pi/2$.

For completeness, choose
\[
 M_0=\max\{\lambda_0,4a/\pi,1\},\quad
 T_0=(M_0+L_0)/a,\quad
 T_1=T_0+\frac2\pi\log(M_0/\lambda_0).
\]
For $t\geq T_0$ while $\lambda\leq M_0$,
$\dot\lambda\geq\pi\lambda/2$, so $\lambda$ reaches $M_0$ by $T_1$.
The line $b(t)=M_0+a(t-T_1)$ satisfies
$b(t)-at=M_0-aT_1\leq-L_0$.
At a contact point $\lambda=b$ for $t\geq T_1$,
$\dot\lambda\geq\pi b/2\geq2a>\dot b$.
Scalar comparison gives $\lambda\geq b$ thereafter. Before $T_1$
we have $\lambda>0$, so $\lambda(t)\geq at-aT_1$ for all $t\geq0$.
\end{proof}

\begin{proof}[Proof of Theorem~\ref{thm:introcancellation}]
Global existence is Lemma~\ref{lem:massvariation}.
The contraction and gap statements follow from
Lemmas~\ref{lem:Xcontract} and \ref{lem:growinggap}.
The latter also gives
$\norm{e^{-B}}=\norm{e^{-L}}\leq e^{aT_1}e^{-at}$.
Thus
\begin{align*}
 \norm{\Omega_V}_1+\norm{\Omega_U}_1
 &\leq\bigl(2n|z_0|e^{aT_1}+n(\kappa+1)s(0)\bigr)e^{-at},\\
 \norm{\Omega_W-z_1\id}_1&\leq s(0)e^{-at}.
\end{align*}
The central-charge identity is $\tr\Omega=m_0z_1$.
The triangle inequality gives the lower mass bound
$m(h(t))\geq m_0|z_1|$ and, with the last two estimates, the upper
bound \eqref{eq:intromass}.

It remains to identify convergence of a metric in fixed coordinates.
For $r\geq0$,
\[
 0\leq\arg(z_1-r)-a
 =\int_0^r\frac{\Ima z_1}{|z_1-u|^2}\,du
 \leq\frac r{\Ima z_1}.
\]
Spectral calculus in \eqref{eq:targetX} yields
$0\leq\Phi_W-a\id\leq s(0)e^{-at}(\Ima z_1)^{-1}\id$.
For $\widehat h_W=e^{at}h_W$, comparison of Hermitian forms gives
\[
 e^{-s(0)/(a\Ima z_1)}h_W(0)
 \leq\widehat h_W(t)\leq h_W(0).
\]
The normalized metric is decreasing, and hence has a positive definite
limit. Its derivative in any fixed coordinate norm is $O(e^{-at})$,
by the same comparison and the phase bound. Integration proves
exponential convergence.
\end{proof}

\begin{remark}\label{rem:zero}
Zero-dimensional cases have direct descriptions. If $V=U=0$, the only
coefficient is $z_1\id_W$ and $h_W(t)=e^{-at}h_W(0)$.
If $W=0$, all arrows vanish; writing $\theta_0=\arg z_0$ gives
$\Phi_V=\theta_0\id$ and $\Phi_U=(\theta_0-\pi)\id$.
Consequently $D(t)=e^{\pi t/2}D(0)$ in moving frames and
$m(h(t))=2|z_0|\tr e^{-e^{\pi t}D(0)^\dagger D(0)}$.
Thus the acyclic mass decays faster than every fixed exponential.
\end{remark}
